% Library unit: opposite-coloured bishops, the second defensive system. Printed once, before the
% first position that uses it; see lib/README.md.

\paragraph{Opposite-coloured bishops: the second defensive system}\label{second-defence-system}\index{Opposite-coloured bishops!Second defensive system}
With bishops of opposite colours, the defender can often stop two passed pawns even without
blockading both with his king. In the second defensive system the king blockades one pawn, while
the bishop alone holds the other pawn from a distance, by controlling a square in front of
it~\cite{Dvoretsky:2008-endgame-manual}.

The attacker breaks it in two steps. First, a king invasion:\index{Opposite-coloured bishops!King invasion}
his king walks round to the pawn the bishop is holding and supports its advance on the squares the
defending bishop cannot control. Second, when that pawn reaches the square the bishop guards, the
defender must give up his bishop for it, and the remaining ending of bishop and pawn against the
king is won.\index{Opposite-coloured bishops!Bishop given up for a pawn} So the defender's king must
also keep watch: it has to stop the attacking king from reaching the pawn that the bishop holds.
